\section{Ce que ce document ne démontre pas}
\label{sec:limites}

Un document qui descend au trade individuel donne une impression de précision
qui peut tromper. Voir trente-cinq lignes détaillées avec des prix au centième
ne rend pas l'échantillon plus grand ni le backtest plus réel. Cette section
énumère ce qu'il faut retenir \emph{contre} les chiffres qui précèdent.

\subsection*{L'échantillon est minuscule}

\textbf{Trente-cinq trades}, dont \textbf{cinq} portent l'intégralité du
résultat. Retirer la seule position d'août~2025 ramène le résultat net de
$+$32\,141\,\$ à $+$10\,490\,\$. Aucune statistique par sous-groupe --- par
année, par régime, par durée --- ne dispose ici de la moindre puissance. Les
ventilations présentées sont \emph{descriptives}~: elles racontent ce qui s'est
passé, elles ne fondent aucune inférence.

\subsection*{Le simulateur a tourné en mode interpolé}

Le backtest utilise le modèle \emph{OHLC 1~minute}, qui reconstitue les
exécutions à l'intérieur des barres au lieu de rejouer les cotations réelles.
Ce mode flatte systématiquement les résultats. \textbf{Tout chiffre
d'exécution présenté ici est donc un majorant}, et non une mesure. Le passage
au mode «~toutes cotations réelles~» suppose de télécharger l'historique de
cotations depuis une session connectée au courtier.

\subsection*{Le repli affiché est une borne inférieure}

Le repli de 46,1\,\% est calculé sur la balance, qui ne bouge qu'à la clôture
des positions. Le capital réellement exposé fluctue entre deux clôtures sans
apparaître sur cette courbe. C'est aussi la raison pour laquelle \textbf{aucun
ratio de Sharpe n'est publié ici pour le moteur or seul}~: une mesure de
dispersion calculée sur une balance est fausse par construction. Les ratios de
Sharpe du dossier proviennent du simulateur lui-même, sur le portefeuille
complet.

\subsection*{Le moteur ne trade pas avant novembre 2021}

Le moteur a besoin de 252~séances d'historique quotidien avant de produire un
signal. Sur ce backtest, l'historique du courtier ne le permet qu'à partir du
\textbf{9~novembre~2021}, alors que la simulation démarre le
1\textsuperscript{er}~janvier~2021. Les dix premiers mois sont donc vides par
construction, et non par décision du moteur. Toute performance annualisée
calculée sur la fenêtre complète est mécaniquement diluée.

\subsection*{Deux postes restent non attribués}

L'attribution du résultat est exacte pour le prix, le portage et la commission.
En revanche, l'effet de l'\textbf{arrondi des volumes} et celui du
\textbf{stop de sécurité} sur la distribution des trades ne sont pas
chiffrables à partir des artefacts disponibles. Ils sont nommés, leur ordre de
grandeur est indiqué quand il est connu, et ils ne sont dilués dans aucun
résidu.

\subsection*{Le passé de l'or n'est pas une promesse}

La fenêtre étudiée couvre une hausse historique de l'or, de 1\,800 à plus de
5\,000~dollars l'once. Un moteur de suivi de tendance long uniquement ne
\emph{peut} que bien s'y comporter. Rien dans ces trente-cinq trades ne
démontre que le moteur produirait un résultat comparable sur une décennie de
marché sans direction --- et les années 2022 et 2023, plates à négatives,
donnent un aperçu de ce à quoi ressemblerait cette période.

\begin{warningbox}
\textbf{Conclusion opérationnelle.} Ce document éclaire le \emph{fonctionnement}
du moteur or~: comment il entre, combien de temps il tient, ce qu'il paie, où
il diverge de sa version de recherche. Il ne modifie en rien la recommandation
du dossier, qui reste un déploiement en simulation avant tout capital réel.
\end{warningbox}
